\section{Related Work}
\label{sec:related_work}

\subsection{Pre-Deep-Learning DTI Methods}

Early DTI methods cover three families that remain conceptually relevant.
Chemogenomic and similarity-inference approaches framed DTI as a link-prediction
problem over drug and target spaces, exploiting the heuristic that chemically
similar drugs bind functionally related proteins
\citep{yamanishi2008prediction,keiser2007relating}. Kernel learning and
matrix-factorization variants formalized this view by recovering latent drug
and target factors from sparse interaction matrices, regularized by
neighborhood or profile similarity \citep{liu2016nrlmf,he2017simboost}.
Heterogeneous-network methods then extended pairwise similarity to broader
biomedical relations spanning drugs, targets, diseases, side effects, and
knowledge graphs \citep{luo2017dtinet,wan2019neodti,ye2021kgenfm}. Classical
machine-learning predictors, including support vector machines, random forests,
nearest-neighbor classifiers, logistic regression, and boosting, served as
strong baselines whenever molecular descriptors, protein features, and curated
negatives were available
\citep{cortes1995svm,breiman2001random,guo2003knn,friedman2000logistic,liu2015credible}.
These methods are interpretable and data-efficient but rely on curated
similarity definitions and coarse pairwise summaries, which limits their
ability to capture residue-level patterns, pair-specific molecular
substructures, or three-dimensional binding geometry.

\subsection{Drug Representation}

Modern small-molecule encoders span the 1D, 2D, and 3D structural views of a
ligand. The one-dimensional view treats a molecule as a token sequence:
SMILES strings provide a compact line notation for molecular graphs
\citep{weininger1988smiles}, molecular fingerprints such as ECFP encode local
substructure patterns into fixed-length vectors used in classical predictors
and similarity search \citep{rogers2010ecfp}, and broader cheminformatics
descriptors summarize physicochemical and topological properties
\citep{todeschini2009descriptors}. Pretrained molecular language models extend
the 1D view through self-supervised SMILES modeling: ChemBERTa, for example,
learns transferable chemistry features from large molecule corpora that benefit
downstream tasks including DTI \citep{chithrananda2020chemberta}.

Two-dimensional graph encoders replace linear notations with explicit
atom-bond topology. Neural graph fingerprints and message-passing networks made
molecular representation learning differentiable and end-to-end
\citep{duvenaud2015graphfps,gilmer2017mpnn}, and DTI-specific graph models such
as GraphDTA, DGraphDTA, and MGraphDTA exploit molecular topology, protein
contact maps, or multiscale graph representations to better capture chemical
neighborhoods and their relation to binding
\citep{nguyen2021graphdta,jiang2020dgraphdta,yang2022mgraphdta}.

Three-dimensional encoders introduce binding-relevant geometry. Continuous-filter
networks such as SchNet model coordinate-dependent molecular signals
\citep{schutt2017schnet}, and equivariant graph neural networks (EGNN) capture
the rotational and translational structure of conformer atoms with low overhead
\citep{satorras2021egnn}. Existing multimodal DTI pipelines that include a 3D
branch typically aggregate candidate conformers as a drug-side property
\citep{satorras2021egnn,ji2026ldmdti}, so the same ligand contributes the same
geometric summary regardless of which protein it is paired with --- a design
choice that we revisit in Section~\ref{sec:methodology}.

\subsection{Protein Representation}

Protein-side modeling has shifted from hand-crafted descriptors to pretrained
sequence encoders. Traditional features such as pseudo-amino-acid composition,
sequence-similarity scores from BLAST-derived profiles, and domain or
functional annotations from Pfam and Gene Ontology provide biologically
interpretable signals, although they tend to miss fine-grained residue
dependencies that are not captured by fixed descriptors
\citep{chou2001pseaac,altschul1997psiblast,finn2014pfam,ashburner2000go,dubchak1995global}.
Deep learning then introduced learned encoders: convolutional and recurrent
networks capture local residue motifs and short-range dependencies
\citep{lee2019deepconv,karimi2019deepaffinity}, while Transformer-style
architectures model longer-range residue interactions through self-attention
\citep{chen2020transformercpi}. Protein language models broadened this
direction by pretraining on large protein corpora: ProtTrans and ProtBERT
learn contextual residue embeddings from self-supervised sequence modeling,
and ESM shows that scaled unsupervised learning recovers structural and
functional signals from sequence alone
\citep{elnaggar2021prottrans,rives2021esm}. Pair-level extensions further use
protein language space as the modeling substrate itself: ConPLex aligns drug
and protein embeddings in a shared contrastive space, and DLM-DTI combines
dual language models with a teacher-student hint mechanism that injects
target-aware signal into a smaller student
\citep{singh2023complex,lee2024dlmdti}. These backbones supply strong
general-purpose target embeddings; they are nevertheless learned independently
of the specific drug partner that will be considered downstream, which leaves
pair-specific protein modulation to later stages of the pipeline.

\subsection{Interaction Modeling}

After per-modality encoders, a DTI model must decide compatibility between drug
and target representations. The simplest design --- concatenation of unimodal
embeddings before an MLP head --- is broadly compatible across encoders but
offloads all cross-modal reasoning onto compressed global vectors. Attention
modules sharpen this by aligning compound substructures with protein residues:
TransformerCPI applies sequence self-attention to compound-protein pairs
\citep{chen2020transformercpi}, MolTrans learns substructure-level interactions
between drug and protein tokens with Transformer and convolutional modules
\citep{huang2021moltrans}, and IIFDTI separates pair-specific (interactive) and
modality-specific (independent) features through an attention design
\citep{cheng2022iifdti}. Bilinear and bidirectional models raise the resolution
of cross-modal matching further: DrugBAN uses bilinear attention to capture
interpretable drug-target associations and add domain adaptation across
distributions \citep{bai2023drugban}, and BINDTI combines a convolution-attention
hybrid (ACmix) protein encoder with bidirectional multi-head attention that
explicitly models drug-to-protein and protein-to-drug intention streams
\citep{peng2024bindti}. Cross-attention frameworks such as MFCADTI integrate
multiple feature views before prediction \citep{quan2025mfcadti}, while
interaction-fingerprint methods rethink how pairwise evidence is represented:
PSICHIC learns hierarchical protein-ligand interaction fingerprints that
incorporate physicochemical channel information
\citep{koh2024psichic}, and BarlowDTI uses a Barlow-Twins-style self-supervised
objective to learn modern one-dimensional DTI representations under stronger
split and representation controls \citep{schuh2025barlowdti}. Attention-based
interaction modules have therefore moved DTI prediction beyond late
concatenation; most current designs still perform pair-specific reasoning only
after independent unimodal encoders have compressed their inputs, however,
which can discard substructure-level or residue-level evidence that a more
deeply coupled design would preserve.

\subsection{Multimodal DTI Systems and Positioning of TAMR-DTI}

Multimodal frameworks combine the representational threads above rather than
choose a single one. LDM-DTI integrates pretrained drug and protein language
features with graph topology, equivariant geometric encoders, and a dynamic
bidirectional interaction module \citep{ji2026ldmdti}. DrugLAMP couples
pretrained biochemical representations with a multi-attention pooling head
designed for transferable DTI prediction across data distributions
\citep{luo2024druglamp}. GEFormerDTA fuses graph-structured atom features with
a Transformer encoder through early-fusion blocks for binding-affinity
prediction \citep{liu2024geformerdta}. EviDTI adds an evidential learning
output head on top of a multimodal DTI backbone to quantify predictive
uncertainty alongside binary prediction \citep{zhao2025evidti}. Beyond pretrained
graph models, the rapid progress of protein structure prediction has further
made predicted three-dimensional protein structure an accessible resource for
downstream target modeling \citep{jumper2021alphafold}, although the present
work focuses on sequence-derived protein representations to remain comparable
with existing BiIntention-style benchmarks.

Most directly related to our protein-side use of state-space modeling, BiMA-DTI
explores a Mamba-Attention hybrid for DTI prediction in which Mamba and
self-attention blocks alternate to model long-range dependencies on both drug
and protein streams \citep{shui2025bimadti}. TAMR-DTI differs from this hybrid
in two ways. First, the protein-side Mamba refinement is applied only to the
ordered protein token stream, so that drug tokens --- which do not carry a
biologically meaningful linear order --- are not subjected to artificial
sequential dynamics. Second, the bidirectional drug-target matching head
(BiIntention) is preserved as a separate component rather than blended with the
state-space block, so that explicit cross-modal matching and long-range
sequence refinement are not entangled in a single fusion stage. In the main
benchmark comparison in Section~\ref{sec:experiments}, we report GraphDTA,
DrugBAN, IIFDTI, BINDTI, TransformerCPI, MolTrans, and LDM-DTI as the closest
published baselines for which protocols, splits, and reported metrics are
directly compatible with our repeated-seed evaluation; the other methods
discussed above use different splits, evaluation protocols, or unavailable
training code, which prevents like-for-like reproduction under the same
protocol and limits them to contextual discussion.
